\documentclass[12pt]{article}
\usepackage[margin=1in]{geometry}
\usepackage{amsmath,amssymb,amsthm}
\usepackage{graphicx}
\usepackage{booktabs}
\usepackage[numbers,sort&compress]{natbib}
\usepackage{hyperref}

\newtheorem{proposition}{Proposition}

\title{Graduated, Restorative Enforcement for Trustless Bitcoin Vaults:\\
Pre-Signed Soft Liquidation with Solvency-Sized Buyback}
\author{Jianming Liu\\DefinanceX Capital}
\date{Companion note to \emph{Target Health Factor based Soft Liquidation with
Collateral Restoration for DeFi Lending} --- working draft}

\begin{document}
\maketitle

\begin{abstract}
Trustless Bitcoin vaults enforce loan terms through pre-signed transactions,
covenant committees, and oracle attestations---and precisely because no party
holds discretion at execution time, every deployed and proposed design we are
aware of enforces liquidation as a \emph{binary} event: one hard price
threshold, full seizure of the vault \cite{babylon-trustless-vaults,firefish-protocol}.
For the BTC-denominated borrower this is the worst possible failure mode: a
single oracle print can permanently remove their entire BTC position. We show
that graduation and restoration are compatible with pre-signing. First,
target-HF partial liquidation compiles to a \emph{geometric price ladder} whose
rung prices and quantities are closed-form at vault creation
(Proposition~\ref{prop:ladder}), hence pre-signable as tranched spending paths.
Second, restoration is where pre-signing bites: the natural pre-signable rule
---buy back exactly what was sold (fixed quantities)---was our original design,
and it suffers severe ping-pong re-liquidation; on BTC 2019--2024 it executes
$45\times$ more transactions for essentially zero net benefit. The cure from
our companion paper, a solvency- and stress-sized buyback quantity
$B^{\star}_{\text{stress}}(C,D,P)$, appears state-dependent and therefore
un-pre-signable, but we show it compiles back to a finite pre-signed structure:
constraining the band to unwind last-sold-first-bought collapses the reachable
states to a contiguous-suffix family of size $O(N^2)$, each with closed-form
$(C,D)$, so safe quantities can be computed and signed at origination
(DLC-style, per oracle price bucket). Empirically, on identical BTC paths at
native cadence, graduation alone halves borrower loss and cuts worst-case bad
debt by three quarters versus full seizure, and the sized restoration matches
the fixed design's outcome with $1/20$ of the executions.
\end{abstract}

\section{The problem: trustlessness forces binary enforcement}
\label{sec:problem}

Native-BTC lending holds collateral in Bitcoin script: multisig escrow with
pre-signed transactions (Firefish \cite{firefish-protocol}), discreet log
contracts (Lava), or, most recently, trustless vaults combining pre-signed
spending conditions, a covenant committee, and oracle attestations
(Babylon \cite{babylon-trustless-vaults}). The defining property of all three
is that \emph{every state transition must be authorized in advance}: at
execution time there is no party with discretion, so there is nothing to
negotiate and no quantity to compute.

The consequence is visible in the deployed designs. Babylon's vault paper
specifies liquidation as: \emph{``if the price oracle releases a signature
certifying that the BTC price has dropped below \$50,000''}, the liquidator
repays the debt and \emph{``can redeem the BTC''}---the entire vault, in one
step \cite{babylon-trustless-vaults}. Firefish prevents this outcome only by
conservatism (50\% LTV and a multi-step margin call), not by graduation.
There is no partial liquidation and no restoration anywhere in the stack:
a single threshold crossing permanently converts the borrower's whole BTC
position. On gap moves this is also bad for the \emph{lender}: the one
liquidation executes far below the threshold, so full seizure coexists with
large bad debt (Table~\ref{tab:native_btc_comparison}).

On EVM lending, graduated liquidation already exists (Aave v4's target-HF
engine) and our companion paper adds a guarded restoration leg as an
incremental, optional module. On native BTC the situation is inverted: the
baseline is the hard-threshold full seizure of 2019-era DeFi, the borrower
population is BTC-denominated by construction, and the binding constraint is
that any improvement must be \emph{pre-committable}. This note shows the two
mechanisms of the companion paper---the target-HF band and the solvency-sized
buyback---are exactly the closed-form objects that pre-signing requires.

\section{Graduation compiles to a pre-signable geometric ladder}
\label{sec:ladder}

Consider the companion paper's sell rule: when $\mathrm{LTV}_t \ge
\mathrm{LLTV}$, repay
$\Delta D^{\star} = D\,(\mathrm{HF}^{\star}-\mathrm{HF}_t)/
(\mathrm{HF}^{\star}-\mathrm{LT}(1+b))$, restoring the health factor to
$\mathrm{HF}^{\star}$. Online, this is a state-dependent rule. Under monotone
decline, however, it is equivalent to a fixed ladder computable at origination.

\begin{proposition}[Pre-computable liquidation ladder]
\label{prop:ladder}
Fix $\mathrm{LT}$, $\mathrm{LLTV}$, $\mathrm{HF}^{\star}>\mathrm{LT}/\mathrm{LLTV}$,
bonus $b$, and an initial position $(C_0, D_0)$ first triggering at price
$P_1$. Under monotone price decline with continuous monitoring, the $k$-th
partial liquidation of the online target-HF rule occurs at
\[
P_k \;=\; P_1\, r^{\,k-1},
\qquad
r \;=\; \frac{\mathrm{LT}}{\mathrm{HF}^{\star}\,\mathrm{LLTV}} \;<\; 1,
\]
and repays the constant debt fraction
\[
\varphi \;=\; \frac{\mathrm{HF}^{\star}-\mathrm{LT}/\mathrm{LLTV}}
{\mathrm{HF}^{\star}-\mathrm{LT}(1+b)},
\qquad
D_k = D_0 (1-\varphi)^k ,
\]
with sale quantities $q_k = (1+b)\,\varphi D_{k-1}/P_k$. All rung prices and
quantities are therefore closed-form at vault creation.
\end{proposition}

\begin{proof}
At every trigger $\mathrm{LTV}=\mathrm{LLTV}$, i.e.\
$\mathrm{HF}=\mathrm{LT}/\mathrm{LLTV}$, so the repair fraction
$\Delta D^{\star}/D = \varphi$ is the same at every rung. The repair moves the
position to $\mathrm{HF}=\mathrm{HF}^{\star}$, i.e.\
$\mathrm{LTV}=\mathrm{LT}/\mathrm{HF}^{\star}$; with debt and collateral then
fixed, $\mathrm{LTV}\propto 1/P$, so the next trigger
($\mathrm{LTV}=\mathrm{LLTV}$) occurs after a further price factor
$(\mathrm{LT}/\mathrm{HF}^{\star})/\mathrm{LLTV} = r$.
\end{proof}

With the baseline parameters ($\mathrm{LT}=0.83$, $\mathrm{LLTV}=0.85$,
$\mathrm{HF}^{\star}=1.05$, $b=0.01$): $r = 0.930$ and $\varphi = 0.347$ ---
rungs every $7\%$ down, each retiring about a third of the remaining debt.

\paragraph{Vault mapping.} Tranche the collateral UTXO at origination so that
rung $k$'s quantity $q_k$ sits in its own output, with a spending path
authorized by an oracle attestation of $P \le P_k$ (plus debt-repayment proof,
exactly as in the Babylon liquidation flow). Discrete oracle cadence makes
execution conservative relative to the continuous idealization; our evaluation
below uses daily cadence throughout. \emph{TODO: dust/fee analysis for the
number of rungs $N$; re-signing checkpoints if the position is modified.}

\section{Restoration: the fixed-buyback trap and its cure}
\label{sec:restoration}

\paragraph{The original design and its disease.} The restoration leg buys
collateral back at the band bottom on recovery, reborrow-funded. The natural
pre-signable quantity rule---and our original design---is \emph{fixed}: buy
back exactly the rung quantities that were sold. It does not require state,
so it can be signed at origination. But it is unsafe: a reborrow-funded buy
of USD value $x$ adds equal collateral and debt and therefore \emph{lowers}
the health factor; buying back a full rung mid-recovery re-arms the trigger,
and the position ping-pongs---sell, buy, re-sell---burning the liquidation
bonus (and, on Bitcoin, transaction fees and spread) on every cycle. The
companion paper measured this on EVM paths (the unguarded leg increases
borrower loss, $p\approx 7\times10^{-8}$); Table~\ref{tab:native_btc_comparison}
measures it at native cadence: $3\times$ the liquidation events and $45\times$
the executions of the sized rule, for essentially zero net benefit.

\paragraph{The cure.} The companion paper's sizing rule caps each buy at
\[
B^{\star}_{\text{stress}}(C,D,P) \;=\;
\max\!\left(0,\;
\frac{\mathrm{LT}(1-d)\,C P - D}{1-\mathrm{LT}(1-d)}\right),
\]
the largest reborrow such that the post-buy health factor remains $\ge 1$ even
after a further fractional drop $d$. This eliminates ping-pong, but it is a
function of the live state $(C,D,P)$---seemingly incompatible with
pre-signing.

\paragraph{Compiling the cure back to pre-signed form.} The state dependence
is an illusion once the unwind order is fixed:

\begin{proposition}[LIFO compilation, sketch]
\label{prop:lifo}
Constrain the band to unwind last-sold-first-bought: rung $j$ may be bought
back only if all rungs $j' > j$ sold after it have been bought back. Then the
reachable inventory states are the contiguous-suffix family
$\{(k,j): 0 \le j \le k \le N\}$ (sold down to rung $k$, bought back down to
rung $j$), of size $O(N^2)$, and each state's $(C_{k,j}, D_{k,j})$ is
closed-form at origination by Proposition~\ref{prop:ladder} plus the fixed
buy levels. Hence $B^{\star}_{\text{stress}}$ can be evaluated per (state,
oracle price bucket) at origination, and the corresponding buyback
transactions pre-signed, DLC-style.
\end{proposition}

\emph{TODO: full proof with rebound-interleaving paths (buys between sells);
transaction-count accounting $O(N^2 \times \text{buckets})$ against DLC CET
practice; what fraction of the online rule's benefit survives bucketing.}

In architectures with on-chain verification (Babylon's covenant committee and
zero-knowledge proofs), compilation is unnecessary: the closed-form inequality
$x \le B^{\star}_{\text{stress}}(C,D,P_{\text{oracle}})$ is trivially provable,
and the quantity rule can be enforced directly.

\section{Evaluation at native cadence}
\label{sec:eval}

We compare four enforcement designs on identical BTC paths (daily Coinbase
prices 2019--2024, eight non-overlapping 270-day windows, three initial CR
levels, healthy starts; the window set includes the March-2020 single-day
crash). Daily cadence approximates the oracle/confirmation tempo of native-BTC
execution. The full-seizure arm models the Babylon flow (liquidator repays the
debt and redeems the vault's BTC).

\IfFileExists{tables/native_btc_comparison_table.tex}{
\begin{table}[h!]
\centering
\caption{Four enforcement designs on identical BTC paths (2019--2024, eight non-overlapping 270-day windows $\times$ three initial CR levels, daily cadence $\approx$ native-BTC oracle tempo). Graduation alone halves borrower loss and cuts worst bad debt by $3/4$ versus full seizure; the fixed buyback of the original pre-signed design churns ($3\times$ the liquidations, $45\times$ the executions) for $\approx$zero net gain; the solvency/stress-sized rule matches its outcome with $1/20$ of the executions.}
\label{tab:native_btc_comparison}
\resizebox{\textwidth}{!}{%
\begin{tabular}{@{}lrrrrr@{}}
\toprule
Design & BTC retention $\rho$ & Mean loss (USD) & Worst bad debt (USD) & Sells/run & Buys/run \\
\midrule
Hard threshold, full seizure (status quo) & 0.375 & 45,303 & 29,808 & 0.62 & 0.00 \\
Pre-committed target-HF ladder & 0.464 & 22,497 & 7,260 & 2.12 & 0.00 \\
Ladder + fixed buyback (original design) & 0.466 & 22,388 & 7,051 & 6.25 & 40.62 \\
Ladder + solvency/stress-sized buyback & 0.466 & 22,482 & 7,260 & 2.12 & 0.92 \\
\bottomrule
\end{tabular}
}
\end{table}

}{}

Three observations. (i) \textbf{Graduation is the headline}: replacing the
hard threshold with the pre-committed ladder halves mean borrower loss
($45{,}303 \to 22{,}497$ USD), raises BTC retention from $0.375$ to $0.465$,
and cuts worst-case bad debt by three quarters ($29{,}808 \to 7{,}260$)---full
seizure is bad for \emph{both} sides, because on a gap the single liquidation
executes far below the threshold. In the March-2020 window the contrast is
starkest ($89{,}860$ vs.\ $67{,}312$ mean loss). (ii) \textbf{The fixed
buyback churns}: $6.25$ sells and $40.6$ buys per run versus $2.1$ and $0.9$
for the sized rule, with net outcomes indistinguishable---in a setting where
every execution is an on-chain transaction, this is pure friction, and it is
the quantitative form of the ping-pong observed in the original design.
(iii) \textbf{The sized rule is execution-minimal}: same retention and loss
as fixed buyback at $1/20$ the executions, and no induced re-liquidation
($2.1$ sells, equal to the no-buyback ladder).

\emph{Honest caveats: eight independent windows, one asset; at daily cadence
within 270-day windows the restoration legs are largely dormant, so this
table quantifies graduation and churn, not restoration upside (see the
companion paper's crash-regime results for the latter); the oracle is trusted
in all arms, as in the deployed designs.}

\section{Related work}
Babylon trustless vaults (hard threshold, full seizure)
\cite{babylon-trustless-vaults}; Firefish and DLC lending (pre-signed,
full liquidation, prevented only by conservative LTV)
\cite{firefish-protocol}; Curve LLAMMA (graduated and restorative, but
AMM-mediated and crvUSD-denominated---not expressible as pre-signed Bitcoin
spending paths) \cite{curve-llamma}; Aave v4 target-HF (graduated, sell-only,
EVM) \cite{aave-v4-liquidations}; reversible call options (averts liquidation
via external top-ups rather than restoring sold collateral)
\cite{mahmood2023miqado}. To our knowledge no prior design offers graduated
\emph{and} restorative enforcement compatible with pre-signed execution.

\section{Conclusion}
In trustless vaults, enforcement has been binary not because anyone prefers
cliffs, but because pre-commitment seemed to demand fixed actions. The
target-HF band turns out to be a geometric, origination-time-computable
ladder, and the solvency-sized restoration---whose fixed-quantity ancestor we
built first and watched ping-pong---compiles back to a finite pre-signed
structure under LIFO unwind. Graduated, restorative, BTC-denominated lending
is therefore implementable on today's trustless-vault architectures, and the
measured stakes are large: half the borrower loss and a quarter of the bad
debt of the status quo, before any restoration upside.

\paragraph{Reproducibility.} The four-arm comparison and its table are
generated by \path{runs/sweeps/native_btc_comparison.py} in the companion
repository; the engine, data, and the companion paper's full evaluation
pipeline are described there.

\bibliographystyle{abbrvnat}
\bibliography{references}

\end{document}
